\chapter{Lenguaje de hipótesis (DSL)}
\label{sec:dsl}

\section{Propósito}

El DSL de GraphHunter describe \emph{hipótesis} de caza: patrones
que buscamos en el grafo. Una hipótesis es una secuencia de
\emph{steps}: entidades conectadas por relaciones, potencialmente
con restricciones temporales y predicados.

\section{Gramática v1 (estable)}

En notación BNF extendida:

\begin{lstlisting}[basicstyle=\ttfamily\small]
hypothesis  ::= step (relation step)* constraints?

step        ::= entityType ( '{' predicates '}' )?

entityType  ::= 'IP' | 'Host' | 'User' | 'Process'
             | 'File' | 'Domain' | 'Registry' | 'URL'
             | 'Service' | '*'
             | customEntityType           (*@\textrm{si hay extensiones}@*)

relation    ::= '-[' relationType ']->'
             | '<-[' relationType ']-'

relationType::= 'Auth' | 'Connect' | 'Execute' | 'Read'
             | 'Write' | 'DNS' | 'Modify' | 'Spawn'
             | 'Delete' | '*'
             | customRelationType         (*@\textrm{si hay extensiones}@*)

predicates  ::= predicate (',' predicate)*
predicate   ::= field op value
              | builtinPredicate '(' args ')'

constraints ::= '{' constraint (',' constraint)* '}'
constraint  ::= 'k'      '=' int         (*@\textrm{k-simplicity cap}@*)
              | 'window' '=' duration    (*@\textrm{ventana temporal}@*)
              | 'order'  '=' 'temporal' | 'any'

duration    ::= int ('s'|'m'|'h'|'d')
\end{lstlisting}

\section{Entity types built-in}

\begin{table}[H]
\centering
\small
\begin{tabular}{@{}ll@{}}
\toprule
Tipo & Semántica \\
\midrule
\code{IP}       & dirección IPv4/IPv6 \\
\code{Host}     & nombre de host o FQDN \\
\code{User}     & principal de identidad (usuario local, AAD, etc.) \\
\code{Process}  & proceso (nombre de ejecutable + hash opcional) \\
\code{File}     & archivo en filesystem \\
\code{Domain}   & nombre de dominio DNS \\
\code{Registry} & clave/valor de registro Windows \\
\code{URL}      & URL \\
\code{Service}  & servicio / aplicación SaaS \\
\code{*}        & comodín: cualquier tipo \\
\bottomrule
\end{tabular}
\caption{Tipos de entidad incluidos en la gramática v1.}
\end{table}

\section{Relation types built-in}

\begin{table}[H]
\centering
\small
\begin{tabular}{@{}ll@{}}
\toprule
Relación & Uso típico \\
\midrule
\code{Auth}    & User $\to$ Host / Service: autenticación \\
\code{Connect} & IP $\to$ IP / Host: conexión de red \\
\code{Execute} & Process $\to$ File / Process: invocación \\
\code{Read}    & Process $\to$ File / Registry: lectura \\
\code{Write}   & Process $\to$ File / Registry: escritura \\
\code{DNS}     & Host $\to$ Domain: resolución \\
\code{Modify}  & Process $\to$ cualquiera: modificación \\
\code{Spawn}   & Process $\to$ Process: padre-hijo \\
\code{Delete}  & Process $\to$ File: borrado \\
\code{*}       & comodín \\
\bottomrule
\end{tabular}
\caption{Tipos de relación en la gramática v1.}
\end{table}

\section{Predicados built-in}

\begin{itemize}
    \item \code{in_subnet(ip, cidr)} — IP en un bloque CIDR.
    \item \code{matches_regex(field, pattern)} — campo contra
    regex PCRE-like.
    \item \code{is_weekend()} — evento ocurrió fuera de horario
    laboral.
    \item \code{within_hours(start, end)} — ventana horaria.
    \item \code{tag(name)} — entidad tiene la etiqueta indicada.
\end{itemize}

Comparadores escalares disponibles: \code{=}, \code{!=}, \code{>},
\code{>=}, \code{<}, \code{<=}. Strings soportan \code{like}.

\section{Ejemplos}

\paragraph{Lateral movement simple.}
\begin{lstlisting}[basicstyle=\ttfamily\small]
User -[Auth]-> Host -[Connect]-> Host {k=3}
\end{lstlisting}
Busca un usuario que autenticó contra un host, y desde ese host se
abrió una conexión a otro host. La cota \code{k=3} limita la
explosión combinatoria (simplicidad de hipergrafo).

\paragraph{Proceso que escribe en directorio sensible.}
\begin{lstlisting}[basicstyle=\ttfamily\small]
Process -[Write]-> File {path like "C:\\Windows\\System32\\*"}
\end{lstlisting}

\paragraph{DNS tunneling sospechoso.}
\begin{lstlisting}[basicstyle=\ttfamily\small]
Host -[DNS]-> Domain {matches_regex(value, "^[a-f0-9]{32}\\.")}
{window=1h}
\end{lstlisting}
Hosts que resuelven muchos dominios con prefijo hex de 32
caracteres dentro de una ventana de 1h.

\paragraph{Cadena de spawning.}
\begin{lstlisting}[basicstyle=\ttfamily\small]
Process {name="winword.exe"}
  -[Spawn]-> Process {name="cmd.exe"}
  -[Spawn]-> Process {name="powershell.exe"}
\end{lstlisting}

\section{Extensibilidad (ADR-005)}

El trait \code{DslExtension} permite registrar:

\begin{itemize}
    \item \textbf{Entity types custom}: \code{Container},
    \code{ServiceAccount}, \code{Cluster}, etc.
    \item \textbf{Relation types custom}.
    \item \textbf{Predicados custom}: implementan
    \code{BuiltinPredicate} trait.
\end{itemize}

Una extensión registra su símbolo y su implementación en el
\code{ExtensionRegistry} al arrancar. La gramática v2 (detrás de
\code{feature = "dsl-v2-experimental"}) incluye operadores
temporales explícitos (\code{before}, \code{within}) y notación
para rangos.

\section{Compatibilidad}

Garantía v1: una hipótesis escrita hoy contra v1 sigue parseando y
comportándose igual en todas las versiones futuras que mantengan
\code{grammar v1}. Los cambios compatibles agregan símbolos o
predicados; los cambios incompatibles promueven a v2.
